\section{Introduction}\label{sec:intro}

The two-dimensional XY model has a massless phase without ordinary
long-range order. Its expected large-scale description is Gaussian, but
the spin is an exponential of an angular variable, not the angular
variable itself. Identifying a height-field limit therefore does not by
itself identify the spin field. At criticality a further difficulty is
that a scale-dependent scalar factor may survive between the microscopic
spin and its continuum normalization.

We prove a critical spin-field limit on a square with fixed boundary
spins. The scalar normalization is the actual magnetization at the
center of that square. This choice cancels the microscopic amplitude
without assuming a power law, a logarithmic correction, or the existence
of an asymptotic amplitude. The remaining spatial normalization is fixed
by the lattice Green function.

\subsection{The physical critical parameter}

We first specify the critical inverse temperature using an infinite-volume
quantity, independently of the field limit. For \(R\in\N\), let
\[
 \Lambda_R=[-R,R]^2\cap\Z^2,\qquad
 E_R=\{\{x,y\}\subset\Lambda_R:|x-y|=1\}.
\]
Every unordered edge occurs once. At inverse temperature \(b>0\), the
free-box law is
\begin{equation}\label{eq:free-spin-law}
 \frac1{Z_{R,b}}
 \exp\!\left\{b\sum_{\{x,y\}\in E_R}\cos(\theta_x-\theta_y)\right\}
 \prod_{x\in\Lambda_R}\frac{d\theta_x}{2\pi},
 \qquad \theta_x\in\R/(2\pi\Z).
\end{equation}
There are no exterior edges or fixed spins in this definition. Write
\begin{equation}\label{eq:critical-parameter}
 \begin{split}
 C_b(r)&=\lim_{\substack{R\to\infty\\R\ge r}}
       \E_{R,b}\cos(\theta_0-\theta_{re_1}),\qquad e_1=(1,0),\\
 m(b)&=\lim_{r\to\infty}-\frac1r\log C_b(r),\qquad
 b_c=\inf\{b>0:m(b)=0\}.
 \end{split}
\end{equation}
The free-box correlation and mass limits exist,
\(0<b_c<\infty\), and \(m(b_c)=0\). These are the usual
phase-transition facts for the cosine convention
\eqref{eq:free-spin-law}; see \cite[Theorem 1]{vEL} and the
normalization comparison in \cite[Proposition~1.2]{C}. No critical exponent is included
in the definition of \(b_c\).

\subsection{The normalized spin field}

For each integer \(n\ge1\), set
\[
 D=(-1,1)^2,\qquad D_n=[-1,1]^2\cap n^{-1}\Z^2,
 \qquad D_n^\circ=D_n\cap D.
\]
Let \(E_n\) contain all unordered nearest-neighbor pairs in \(D_n\).
Set \(\theta_x=0\) on \(D_n\setminus D_n^\circ\), and integrate only
the interior angles in the probability law
\begin{equation}\label{eq:pinned-spin-law}
 \frac1{Z_n^0}
 \exp\!\left\{b_c\sum_{\{x,y\}\in E_n}
                         \cos(\theta_x-\theta_y)\right\}
 \prod_{x\in D_n^\circ}\frac{d\theta_x}{2\pi}.
\end{equation}
Its expectation is denoted by \(\E_n^0\). In particular edges incident
to the fixed boundary are retained. Define
\begin{equation}\label{eq:center-magnetization}
 a_n=\E_n^0\cos\theta_0=\E_n^0e^{i\theta_0}>0.
\end{equation}
Angle reflection proves the equality, and the positive Fourier-current
expansion proves strict positivity. Thus \(a_n\) is a specified
observable of \eqref{eq:pinned-spin-law}, not a sequence chosen to fit
a limit.

On functions zero at the boundary let
\[
 (L_nu)(x)=\sum_{\substack{y\in D_n\\|x-y|=1/n}}
                   (u(x)-u(y)),\qquad G_n=L_n^{-1}.
\]
The Laplacian is unscaled: it contains no factor \(n^2\).
Put \(K_*=2/\pi\), and define the complex random distribution
\begin{equation}\label{eq:spin-field}
 S_n=\frac1{n^2}\sum_{x\in D_n^\circ}
 a_n^{-1}\exp\!\left\{\frac{G_n(x,x)-G_n(0,0)}{2K_*}\right\}
 e^{i\theta_x}\delta_x,
 \qquad \delta_x(f)=f(x).
\end{equation}
The center factor and the Green correction in this definition are
the same for every test function and every correlation order.

Let \(\Phi_D\) be the real zero-Dirichlet Gaussian free field with
covariance \((-\Delta_D)^{-1}\). Its kernel has the convention
\[
 G_D(z,w)=\frac1{2\pi}\log\frac1{|z-w|}+R_D(z,w),
\]
where \(R_D\) is smooth across the interior diagonal. For circle
averages \(\Phi_{D,\epsilon}\), taken where the circles lie in \(D\),
let
\begin{equation}\label{eq:target-field}
 V_*=\lim_{\epsilon\downarrow0}
 \exp\!\left\{\frac{i\Phi_{D,\epsilon}}{\sqrt{K_*}}
       +\frac{\E\Phi_{D,\epsilon}^2}{2K_*}\right\}.
\end{equation}
The limit is understood as a complex random distribution tested
against compactly supported smooth functions. This is the
full-variance imaginary Gaussian chaos: its counterterm uses the
entire variance, including its finite spatial part.
In this convention
\[
 \E V_*(f)=\int_D f(z)\,dz,\qquad
 \E\bigl[V_*(z)\overline{V_*(w)}\bigr]
       =\exp\{G_D(z,w)/K_*\},
\]
where the second expression denotes the off-diagonal moment kernel.
Its singularity has exponent \(1/4\). Existence and the precise
mixed-moment interpretation are recalled in
Lemma~\ref{fields:circle-chaos}; see also \cite{JSW2020,B}.

\begin{theorem}\label{thm:main}
For the fixed-boundary critical XY law \eqref{eq:pinned-spin-law},
the distributions \(S_n\) in \eqref{eq:spin-field} converge in law
along the full sequence \(n\to\infty\) to \(V_*\) in
\(H^{-3}_{\mathrm{loc}}(D)\).
Moreover, for every \(k\ge1\), every
\(f_1,\ldots,f_k\in C_c^\infty(D;\C)\), and every collection of
nonnegative integers \(p_1,q_1,\ldots,p_k,q_k\),
\begin{equation}\label{eq:main-moment-convergence}
 \E_n^0\prod_{j=1}^k
       S_n(f_j)^{p_j}\overline{S_n(f_j)}^{\,q_j}
 \longrightarrow
 \E\prod_{j=1}^k
       V_*(f_j)^{p_j}\overline{V_*(f_j)}^{\,q_j}.
\end{equation}
The joint laws of these complex smeared fields converge as well.
\end{theorem}

Here \(H^{-3}_{\mathrm{loc}}(D)\) consists of distributions \(T\)
such that \(\chi T\), extended by zero, belongs to
\(H^{-3}(\R^2)\) for every \(\chi\in C_c^\infty(D)\), with the
usual local topology. The theorem includes nonneutral moments:
there is no requirement that the total number of fields equal
the total number of conjugate fields.

The critical input is the Bessel-height coefficient and the height
and pin limit theorems of \cite{C}. Their exact statements and
conventions are recorded in Section~\ref{sec:heights}; the critical
coefficient there is \(a(b_c)=8\pi\), giving
\(K_*=a(b_c)/(4\pi^2)\). We also use the finite comparisons and field
lemmas identified in \cite{B}. We do not deduce the result by
continuing a low-temperature spin theorem to \(b_c\).
Nor does Theorem~\ref{thm:main} assert an asymptotic formula for
\(a_n\), or by itself transfer the result to the free-boundary
infinite-volume correlation. Such a transfer would require
additional amplitude and boundary-to-bulk information.

\subsection{Historical context and related work}

The distinction between absence of magnetization and slowly decaying
correlations is central to the two-dimensional XY model. Mermin and
Wagner~\cite{MW} established the low-dimensional obstruction to
conventional order for isotropic Heisenberg systems. For classical
planar spins, McBryan and Spencer~\cite{MS} obtained quantitative
power-law upper bounds on correlations. Berezinskii~\cite{Berezinskii}
developed the algebraic low-temperature picture and argued for a
phase transition. Kosterlitz and Thouless~\cite{BKT1973} proposed
the transition mechanism through the behavior of topological defects.
The ensuing renormalization-group theory predicts
the critical exponent \(1/4\) and renormalized stiffness
\(2/\pi\)~\cite{Kosterlitz1974,NK}.

Fr\"ohlich and Spencer~\cite{FS} proved a low-temperature phase with
power-law lower bounds for the rotator model and related systems.
Van Engelenburg and Lis~\cite[Theorem 1]{vEL} established a sharp
exponential-versus-polynomial dichotomy for free infinite-volume XY
correlations, including a polynomial lower bound at the threshold.
Lammers~\cite[Theorem 1]{Lammers2023} related the spin and dual-height
correlation lengths and identified the coincidence of the XY
transition with dual-height delocalization. These results identify
the physical threshold in Equation~\eqref{eq:critical-parameter};
they do not supply the critical spin-field limit.

Gaussian descriptions of angular fluctuations also have rigorous
precedents. Newman and Wu~\cite[Theorem 2]{NewmanWu2018} proved
Gaussian limits for angular-gradient observables in a joint
zero-temperature and continuum regime. The companion
\cite[Theorem 7.6]{B} treats the exponential spin field at sufficiently
low fixed temperature. The critical coefficient and fixed-geometry
Bessel-height inputs used here come instead from
\cite[Section 6 and Theorems 2.4, 3.1, 3.2]{C}.
Our task is to pass from these critical inputs to the spin field,
retaining the actual center magnetization as normalizer and treating
all signed correlations, including nonneutral ones.

The proof begins with integer-current and planar-height
duality~\cite[Section 2]{vEL}. The finite comparisons descend from
the lattice-Gaussian inequalities of Regev and
Stephens-Davidowitz~\cite{RSD}, in their Bessel-height form
in~\cite{C}. The target belongs to the theory of imaginary Gaussian
multiplicative chaos developed by Junnila, Saksman and
Webb~\cite{JSW2020}. To pass from correlations to field laws and
moments, we use the Lee--Yang theorem for planar
rotators~\cite{DunlopNewman1975} and its moment-generating-function
inequalities~\cite{Newman1975}, in the finite-volume formulation
of Newman and Wu~\cite{NewmanWu2019}. The sharp comparison for holes
of positive radius and the control of arbitrarily many annular
interfaces are developed below.

\subsection{Proof strategy}

Fourier expansion turns a product of spins into a ratio of height
partition functions. A spin insertion prescribes an integral winding
of the dual height around its location. We call a height law with
such prescribed windings a magnetic sector. Fixed spin boundary
conditions give free dual heights, modulo one common integer
translation.

The height limit does not directly determine a magnetic-sector
ratio when a puncture has the size of one lattice spacing. We first
keep every puncture at a positive macroscopic radius. For such
fixed domains, Section~\ref{sec:sectors} proves the sharp Gaussian
variational formula. Its upper bound uses divergence-free tests.
Its lower bound introduces quadratic penalties on box averages,
uses separated-pin estimates to compare pieces, and refines the
observation cells to recover the Dirichlet energy. All height limits
occur at fixed geometry.

To return to microscopic punctures, we surround each spin insertion
by a sequence of annuli. At an interface, the gluing kernel is the
partition-function factor for restoring the deleted bonds, conditional
on all height differences in a separating band. The integer constant
between two pieces is still summed. This relative-data convention
is essential both for exact factorization and for the treatment of
the free constant mode.

Sections~\ref{sec:annuli} and~\ref{sec:kernels} give two different
controls on these kernels. After division by their centered
zero-data value, they lie between zero and one for every datum.
Their average loss from one is small on long, approximately circular
annuli, including in a unit magnetic sector. These estimates are
not multiplied into an error proportional to the number of scales.
Instead, a contraction on mean-zero functions shows that an
arbitrary fixed microscopic core eventually produces a last-band
density bounded by two relative to the free annulus law. The core
may have an arbitrarily poor initial gluing factor; positivity is
enough.

We then attach all the annular sequences to one exterior region.
The exterior has positive-radius holes and is governed by the
variational formula already proved. The entire microscopic
contribution is the same for every unit insertion and for the
single insertion defining \(a_n\). Subtracting the corresponding
logarithms cancels that contribution exactly. Section~\ref{sec:correlations}
computes the remaining Green-function interaction and proves
uniform convergence on separated compact configurations.

The final step needs more than separated correlations. A weaker,
integrable two-point bound controls collisions uniformly in \(n\).
Together with positivity and the Lee--Yang moment estimate, it
removes every collision region from every mixed moment, without
requiring a sharp pointwise bound for higher correlations.
Section~\ref{sec:fields} then identifies the moments, proves
moment determinacy, and obtains local Sobolev tightness.

The sharp fixed-hole comparison and the positive-kernel stability
argument are the transferable parts of the proof. They separate
fixed-geometry Gaussian information from the microscopic
normalization, and use averaged control of conditional data where
a uniform conditional limit is unavailable.
